%% ============================================================
%% Pucallpazo — LARR submission
%% Overleaf compatible: pdflatex + biblatex (chicago author-year)
%% ============================================================
\documentclass[12pt]{article}

%% --- Encoding & fonts ---
\usepackage[T1]{fontenc}
\usepackage[utf8]{inputenc}
\usepackage{mathptmx}          % Times New Roman equivalent

%% --- Page layout ---
\usepackage[
  top=1in, bottom=1in,
  left=1in, right=1in
]{geometry}
\usepackage{setspace}
\doublespacing

%% --- Paragraph style (LARR: first-line indent, no extra space) ---
\setlength{\parindent}{0.5in}
\setlength{\parskip}{0pt}

%% --- Tables ---
\usepackage{booktabs}
\usepackage{tabularx}
\usepackage{array}
\usepackage{multirow}

%% --- Hyperlinks (Overleaf-friendly, no colored boxes) ---
\usepackage[
  hidelinks,
  pdfauthor={Anonymous},
  pdftitle={Subnational Critical Junctures from the Periphery}
]{hyperref}

%% --- Bibliography: biblatex with Chicago author-year ---
\usepackage[
  backend=biber,
  style=chicago-authordate,
  autocite=inline
]{biblatex}
\addbibresource{references_zotero.bib}   % <-- replace with your .bib filename

%% --- Footnotes (LARR uses footnotes, not endnotes) ---
\usepackage[bottom]{footmisc}

%% --- Misc ---
\usepackage{microtype}
\usepackage{csquotes}

%% ============================================================
\begin{document}

%% ---- Title block (double-blind: no author names) -----------
\begin{center}
  {\large\bfseries Subnational Critical Junctures from the Periphery:\\
  The Pucallpazo and Territorial Reconfiguration in Peru (1978--1981)}
\end{center}

\vspace{1em}

%% ---- Abstract ----------------------------------------------
\noindent\textbf{Abstract}

\noindent This article examines the Pucallpazo, a cycle of collective
mobilization that took place in the Amazonian city of Pucallpa between
1978 and 1981 and culminated in the creation of the Department of
Ucayali, permanently altering the territorial organization of the
Peruvian state. The case raises a broader question for the study of
institutional change in Latin America: under what conditions can
mobilization from a peripheral region generate durable transformations
in state structure?

Building on historical institutionalism, the article distinguishes
between gradual institutional change and critical junctures and evaluates
whether the creation of Ucayali resulted from incremental adaptation or
from a discontinuous reconfiguration of territorial authority. Using
process tracing based on archival sources, press reports, official
documents, and oral histories, the analysis reconstructs the sequence
linking long-term center--periphery tensions, a permissive national
political context, and the self-organization of a heterogeneous local
coalition.

The findings show that the reform cannot be explained as gradual change.
Instead, the Pucallpazo represents a subnational critical juncture in
which sustained mobilization redefined territorial authority and generated
effects across multiple institutional domains. In comparative perspective,
the case differs from both elite-driven reforms and protest cycles that
yield only temporary concessions. The article argues that critical
junctures may originate outside national elites when structural tensions,
political uncertainty, and coordinated collective action converge,
highlighting the importance of peripheral conflicts for understanding
democratization and state formation in Latin America.

\vspace{0.5em}
\noindent\textbf{Keywords:} critical junctures; subnational
democratization; institutional change; Peru; bottom-up mobilization.

\newpage

%% ============================================================
\section{Introduction}
%% ============================================================

In 1978, a cycle of collective mobilization known as the Pucallpazo
erupted in the Amazonian city of Pucallpa, then part of the Department
of Loreto. What began as a local protest evolved into a sustained
regional movement that, within three years, led to the legal creation
of the Department of Ucayali and a reorganization of Peru's territorial
state. This remains the only instance in Peruvian history in which a
subnational mobilization succeeded in creating a new department. The
case raises a broader question for the study of institutional change in
Latin America: under what conditions can bottom-up mobilization from a
peripheral region produce durable transformations in state structure?

Research on institutional change has traditionally emphasized either
elite negotiation or gradual adaptation. Historical institutionalism
has shown that institutions often evolve incrementally through processes
such as layering, drift, conversion, and displacement rather than through
abrupt rupture \autocite{mahoney_explaining_2010,streeck_introduction_2005}. At the
same time, the literature on critical junctures highlights moments of
heightened contingency in which political choices produce long-lasting
institutional effects \autocite{collier_shaping_1991,capoccia_study_2007}.
These approaches have been highly influential in the study of regime
change, democratization, and state formation in Latin America, where
transitions are often analyzed as outcomes of elite bargaining under
conditions of crisis \autocite{huntington_tercera_2002,odonnell_transitions_1986}.

Parallel to this tradition, a growing body of work has emphasized the
role of collective action and subnational politics in shaping
institutional outcomes. Studies of contentious politics show that
sustained mobilization can alter authority relations even in highly
centralized states \autocite{tilly_social_2004,abers_inventing_2000}. Research on
subnational regimes demonstrates that territorial conflicts frequently
produce institutional innovations that cannot be explained solely by
national-level dynamics \autocite{gibson_boundary_2005,giraudy_politics_2010}. Yet most of
this literature focuses on policy change, participation, or regime
variation, rather than on transformations in the territorial
organization of the state itself.

This article examines the Pucallpazo as a case of institutional change
generated from the periphery. Between 1978 and 1981, a heterogeneous
coalition of local actors---workers, merchants, professionals, regional
authorities, and civic organizations---organized a series of strikes,
demonstrations, and negotiations demanding administrative autonomy from
Loreto. These events took place during the transition from military rule
to civilian government, a period characterized by political uncertainty
and reduced state capacity. The outcome was not merely the expansion of
participation or the adoption of new policies, but the creation of a new
territorial unit with its own administrative and political institutions.

The argument developed here is that the creation of Ucayali cannot be
adequately explained as gradual institutional change. Instead, the case
represents a subnational critical juncture in which long-standing
center--periphery tensions, a permissive national political context,
and sustained collective mobilization converged to produce a durable
redefinition of territorial authority. By reconstructing the sequence
linking structural conditions, political contingency, and collective
action, the article shows how institutional rupture can originate outside
national elites and emerge from conflicts in peripheral regions.

The analysis proceeds as follows. The next section develops the
theoretical framework, distinguishing between gradual institutional
change and critical junctures and specifying the conditions under which
bottom-up mobilization can generate institutional rupture. The third
section describes the research design and sources. Sections four through
six reconstruct the historical background, the cycle of mobilization,
and the institutional outcome. Section seven places the case in
comparative perspective, and the conclusion discusses implications for
the study of democratization, territorial politics, and institutional
change in Latin America.

%% ============================================================
\section{Institutional Change, Critical Junctures, and\\
Bottom-Up Democratization}
%% ============================================================

Institutional transformation can follow multiple paths. In some cases,
institutions evolve gradually as actors reinterpret rules, adjust
enforcement, and modify authority relations over time. In other cases,
political change occurs through relatively short periods of heightened
contingency in which alternative institutional arrangements become
possible and choices made during these moments produce long-lasting
consequences. Distinguishing between gradual change and discontinuous
transformation has been a central concern of historical institutionalism
and comparative politics
\autocite{streeck_introduction_2005,mahoney_explaining_2010,capoccia_study_2007}.

\subsection*{Gradual institutional change}

Research on institutional change has shown that important transformations
often occur without formal rupture. \textcite{streeck_introduction_2005} identify
several mechanisms through which institutions evolve incrementally,
including displacement, layering, drift, conversion, and exhaustion.
These mechanisms allow existing institutional frameworks to adapt to new
conditions while preserving formal continuity.
\textcite{mahoney_explaining_2010} further demonstrate that institutional
stability frequently coexists with gradual shifts in power relations, as
actors reinterpret rules or exploit ambiguities to advance new objectives.

These perspectives have been particularly influential in the study of
Latin American political development, where many reforms have taken place
through negotiated adjustment rather than abrupt breakdown. Democratization
processes, for example, have often been described as the outcome of elite
bargaining under conditions of uncertainty, rather than as the result of
revolutionary rupture
\autocite{huntington_tercera_2002,odonnell_transitions_1986}. From this view,
institutional change is typically cumulative, path-dependent, and
constrained by existing structures.

\subsection*{Critical junctures and discontinuous change}

Alongside gradualist accounts, another tradition emphasizes moments of
discontinuity. The literature on critical junctures argues that
institutional trajectories can be redirected during relatively brief
periods in which structural constraints weaken and political actors face
multiple viable alternatives
\autocite{collier_shaping_1991,capoccia_study_2007}. Decisions made
during such periods may produce long-term consequences that persist even
after the conditions that enabled them disappear.

Critical junctures are not defined solely by the magnitude of change,
but by the combination of contingency, choice, and durability.
\textcite{capoccia_study_2007} stress that a critical juncture requires
the presence of plausible alternative outcomes and the possibility that
different decisions could have led to different institutional paths.
Similarly, \textcite{slater_informative_2010} argue that critical junctures
involve the interaction of antecedent conditions, permissive conditions
that open a window of opportunity, and productive conditions that
determine which of the available alternatives is selected.

This framework has been widely used to analyze regime change, state
formation, and major policy reforms. However, most applications focus on
national-level processes---such as transitions to democracy,
constitutional redesign, or elite pacts---leaving less explored the
possibility that critical junctures may originate in subnational arenas.

\subsection*{Bottom-up mobilization and territorial conflict}

A parallel line of research highlights the role of collective action and
territorial conflict in shaping institutional outcomes. Studies of
contentious politics show that sustained mobilization can alter authority
relations even in highly centralized states
\autocite{tilly_social_2004,abers_inventing_2000}. In such cases, institutional change
does not result from elite design alone, but from pressure generated by
organized actors outside the central government.

Research on subnational politics further demonstrates that territorial
dynamics can produce institutional variation that cannot be explained
solely by national-level factors. \textcite{gibson_boundary_2005} shows that
conflicts over territorial authority often shape the distribution of
power within democratic regimes, while \textcite{giraudy_politics_2010} documents
how subnational political structures may persist or change independently
of national transitions. These findings suggest that institutional change
may originate in peripheral regions where tensions between economic
importance, demographic growth, and limited political representation
become difficult to sustain within existing arrangements.

Despite these insights, the literature has paid limited attention to
cases in which bottom-up mobilization leads not only to policy change or
participation, but to a redefinition of the territorial organization of
the state. The creation of new administrative units, regions, or
provinces represents a particularly demanding form of institutional
change, because it alters formal authority relations and redistributes
resources across levels of government.

\subsection*{Analytical framework}

This article integrates insights from historical institutionalism and
the study of contentious politics to examine whether the creation of the
Department of Ucayali resulted from gradual institutional adaptation or
from a critical juncture generated from below. Following the literature
on gradual change, the analysis first evaluates whether the outcome can
be explained through mechanisms such as layering, conversion, or drift
\autocite{streeck_introduction_2005,mahoney_explaining_2010}. If these mechanisms
cannot account for the transformation, the analysis then assesses whether
the sequence satisfies the criteria for a critical juncture, including
the presence of antecedent tensions, a permissive political context, a
productive mechanism linking mobilization to institutional change, and
durable consequences
\autocite{collier_shaping_1991,capoccia_study_2007,slater_informative_2010}.

By applying this framework to the Pucallpazo, the article seeks to show
that institutional rupture can originate in peripheral regions when
long-term territorial tensions coincide with political uncertainty and
sustained collective action. In such cases, bottom-up mobilization may
not only influence policy, but also redefine the structure of the state
itself.

%% ============================================================
\section{Research Design and Sources}
%% ============================================================

This article employs a qualitative single-case design to evaluate
whether the Pucallpazo constitutes a subnational critical juncture.
Single-case studies do not aim at statistical generalization, but they
are particularly appropriate for theory testing, identification of causal
mechanisms, and reconstruction of contingent historical sequences
\autocite{george_case_2005,bennett_elman_2006}. These features make
them well suited for distinguishing between gradual institutional change
and discontinuous transformation, a distinction that lies at the center
of the present analysis.

Within historical institutionalism, process tracing has become a
standard strategy for assessing causal claims in temporally ordered
sequences \autocite{mahoney_historical_2009,capoccia_study_2007}. Rather than
relying on correlations across cases, process tracing examines whether
the empirical evidence is consistent with the theoretical mechanisms that
a given explanation requires. This approach is particularly useful when
the goal is to determine whether an institutional outcome resulted from
incremental adaptation or from a critical juncture characterized by
contingency and durable effects.

\subsection*{Case selection}

The Pucallpazo is selected as a theoretically relevant case for examining
bottom-up institutional rupture in a highly centralized state. In the
Peruvian context, the creation of the Department of Ucayali represents
the only known instance in which sustained mobilization at the subnational
level led to the formation of a new territorial unit. This makes the case
especially suitable for evaluating whether collective action from the
periphery can produce changes in the formal organization of the state.

The case also presents conditions that are theoretically favorable for
rupture: a long-standing center--periphery cleavage, repeated conflicts
over administrative autonomy, and a period of political uncertainty
associated with the transition from military rule to civilian government.
Because these conditions make institutional change plausible, the case
constitutes a demanding test for gradualist explanations. If the outcome
cannot be fully accounted for by mechanisms of incremental adaptation,
the hypothesis of a critical juncture gains credibility.

\subsection*{Analytical strategy}

The analysis follows a two-step diagnostic strategy derived from the
literature on institutional change. The first step evaluates whether the
creation of Ucayali can be explained through mechanisms of gradual
transformation, such as layering, drift, conversion, or displacement
\autocite{streeck_introduction_2005,mahoney_explaining_2010}. These mechanisms
imply continuity in the basic institutional framework even when important
adjustments occur.

If gradual change cannot fully account for the observed transformation,
the second step assesses whether the sequence satisfies the criteria for
a critical juncture. Following \textcite{collier_shaping_1991} and
\textcite{capoccia_study_2007}, a critical juncture requires a period
of heightened contingency in which multiple outcomes are possible and the
decisions made during that period have lasting consequences.
\textcite{slater_informative_2010} further emphasize the importance of
distinguishing between antecedent conditions, permissive conditions that
open the window for change, and productive conditions that determine
which outcome prevails.

In addition, the analysis considers the magnitude and scope of
institutional change. Transformations that qualify as critical junctures
should be relatively swift, significant, and encompassing in their
effects, altering the distribution of authority or the structure of
political institutions rather than merely adjusting existing rules.

\subsection*{Process tracing}

To evaluate these criteria, the article applies within-case process
tracing \autocite{george_case_2005}. The analysis examines the
sequence linking long-term territorial tensions, the cycle of
mobilization between 1978 and 1981, and the legal creation of the
Department of Ucayali. Evidence is assessed using standard
process-tracing tests, including necessary-condition tests and strong
confirmatory tests, in order to determine whether the observed events
are consistent with the causal mechanisms implied by each explanation
\autocite{bennett_process_2010}.

\subsection*{Sources}

The empirical reconstruction draws on multiple types of sources in order
to triangulate evidence and reduce the risk of retrospective bias. These
include archival documents, official records, press coverage, audiovisual
materials, and oral histories.

Archival sources consist of national and regional government documents
from the late 1970s and early 1980s, including decrees, congressional
debates, and administrative correspondence related to the creation of
the Department of Ucayali. Press coverage from newspapers published in
Lima and in the Amazon region provides contemporaneous accounts of the
mobilization and of the negotiations between local actors and the central
government.

The study also incorporates oral histories and interviews with
participants, journalists, and local authorities. Following the
methodological literature on oral history, testimonies are treated not
only as sources of factual information but also as evidence of how actors
interpreted events and justified their actions \autocite{thompson_voice_2000}.
Whenever possible, these accounts are cross-checked against written
records to ensure consistency.

%% ============================================================
\section{Center--Periphery Tensions in the Peruvian Amazon}
%% ============================================================

The Pucallpazo cannot be understood without considering the long-standing
tensions between the Peruvian state and its Amazonian periphery.
Throughout the twentieth century, the Amazon region occupied an ambiguous
position within the national political economy. Despite its vast territory
and strategic importance, the region remained weakly integrated into
national administrative structures and was often governed through
centralized arrangements that limited local autonomy. Historical accounts
of the Peruvian Amazon repeatedly note the persistence of conflicts over
territorial authority, resource control, and political representation
\autocite{san_roman_perfiles_1994,dourojeanni_amazonia_1990}.

During most of the republican period, administrative organization in the
Amazon favored large territorial units governed from distant capitals.
The Department of Loreto, for example, encompassed extensive areas with
limited infrastructure and difficult communication with Lima. This
administrative structure reinforced dependence on the central government
and reduced the capacity of local actors to influence political decisions.
Studies of regional history emphasize that demands for decentralization
and administrative reorganization appeared repeatedly throughout the
twentieth century, particularly in areas experiencing demographic growth
and economic expansion \autocite{espinoza_soriano_loreto_2016}.

Pucallpa emerged as one of the most dynamic urban centers in the Amazon
during the mid-twentieth century. Its growth was associated with river
trade, timber extraction, and later road connections linking the region
to the central highlands. As the city expanded, local elites, workers'
organizations, and civic associations increasingly questioned the
concentration of political authority in Iquitos, the capital of Loreto.
These tensions reflected a broader pattern in which economic development
in peripheral regions was not accompanied by proportional political
representation \autocite{vivanco_pimentel_gran_1997}.

Demographic change intensified these pressures. Census data from the
late twentieth century show rapid population growth in the Ucayali area
compared to other parts of the Amazon, reinforcing the perception among
local actors that the existing administrative structure no longer
reflected social and economic realities
\autocite{inei_censos_1981,inei_resultados_2018}. Local leaders argued that decisions
affecting the region were taken by authorities who lacked direct
knowledge of local conditions, and that administrative distance
translated into unequal distribution of public resources.

Conflicts over territorial authority were not unique to the Peruvian
Amazon. Comparative research on state formation shows that peripheral
regions often become sites of institutional tension when economic
expansion, demographic change, and limited political representation
coincide \autocite{boone_political_2003}. In such contexts, demands for new
administrative units, regional autonomy, or decentralization may emerge
as attempts to realign political institutions with changing social
conditions.

In Peru, these tensions became particularly visible during the final
years of military rule. The government of General Francisco Morales
Bermúdez (1975--1980) faced growing political opposition while attempting
to manage economic crisis and prepare a transition to civilian government.
This period of uncertainty weakened the capacity of the central state to
control regional conflicts and opened space for local actors to press
demands that had previously been contained.

By the late 1970s, the relationship between Pucallpa and the authorities
of Loreto had become increasingly conflictive. Local organizations accused
the departmental government of neglect, unequal investment, and political
marginalization. Previous petitions for administrative autonomy had
produced limited concessions but no structural change. As a result,
demands for the creation of a new department gradually replaced earlier
requests for minor administrative reforms.

These long-term tensions formed the background against which the
Pucallpazo unfolded. They did not make institutional rupture inevitable,
but they created the conditions under which mobilization could be
interpreted as a struggle over territorial authority rather than as a
temporary protest. When political uncertainty at the national level
reduced the costs of confrontation, these accumulated grievances provided
the basis for sustained collective action that would eventually challenge
the existing administrative order.

%% ============================================================
\section{The Pucallpazo Cycle of Mobilization (1978--1981)}
%% ============================================================

The cycle of mobilization known as the Pucallpazo unfolded between 1978
and 1981 and involved a sequence of strikes, demonstrations, negotiations,
and political confrontations that culminated in the creation of the
Department of Ucayali. Although the demand for administrative autonomy
had precedents, the events of these years differed in scale, coordination,
and political impact. What began as a local protest developed into a
sustained regional movement capable of forcing institutional change.

The first major mobilization occurred in 1978, when local organizations
in Pucallpa called for a general strike demanding greater administrative
autonomy and increased public investment. The protest brought together a
heterogeneous coalition that included labor unions, merchants, transport
workers, professionals, students, and civic associations. Testimonies
from participants emphasize that the movement lacked a single centralized
leadership and instead operated through coordination among multiple
organizations, each with its own agenda but united by the demand for
regional autonomy \autocite{thompson_voice_2000}.

Local sources describe the mobilization as a response to what
participants perceived as systematic neglect by the departmental
authorities in Iquitos. According to regional accounts, public
investment, administrative attention, and political representation were
concentrated in the departmental capital, while rapidly growing areas
such as Pucallpa received limited resources despite their economic
importance \autocite{vivanco_pimentel_gran_1997}. These grievances had
accumulated over decades, but the political context of the late 1970s
allowed them to be expressed in more confrontational forms.

The first Pucallpazo demonstrated the capacity of local actors to
organize large-scale collective action, but it did not immediately
produce institutional change. Instead, it initiated a cycle of
mobilization in which protest, negotiation, and renewed conflict followed
one another over several years. Oral testimonies and journalistic
accounts indicate that the movement persisted because earlier concessions
were perceived as insufficient or unreliable, reinforcing the belief that
only the creation of a new department would guarantee political autonomy.

The second phase of mobilization took place in the context of the
transition from military rule to civilian government. As the regime of
Morales Bermúdez prepared elections and reduced repression, regional
actors perceived an opportunity to press demands that might have been
suppressed under more stable conditions. Interviews with participants
suggest that the expectation of political change at the national level
encouraged local organizations to intensify their efforts, believing
that the central government would be more responsive during a period of
uncertainty.

During this period, the demand for administrative reform became more
clearly defined as a demand for the creation of the Department of
Ucayali. Civic committees and regional leaders coordinated petitions,
public demonstrations, and negotiations with national authorities.
Archival documents and press reports indicate that the issue reached the
national political agenda, forcing government officials to consider
alternatives that had previously been rejected.

In 1980, the return to civilian rule under President Fernando Belaúnde
Terry created a new political context. The incoming government faced
multiple regional demands and limited administrative capacity, which
increased the incentives to resolve conflicts through institutional
concessions rather than repression. According to contemporaneous accounts,
local leaders in Pucallpa interpreted the transition as a window of
opportunity and intensified mobilization in order to secure a definitive
decision.

The creation of the Department of Ucayali was formalized through legal
measures adopted during this period, but the outcome remained uncertain
for several years. Evidence from interviews and official records indicates
that the reform was contested, delayed, and at times threatened with
reversal. Even after the formal decision had been announced, local
organizations continued to mobilize in order to ensure its
implementation.

Audiovisual material, regional chronicles, and interviews with
participants consistently describe the process as one marked by
uncertainty rather than inevitability. Actors involved in the movement
recall that different outcomes were considered possible, including partial
concessions, administrative reorganization without autonomy, or the
cancellation of the reform. The persistence of mobilization played a
decisive role in preventing these alternatives and in sustaining pressure
on the central government until the new department was effectively
established.

%% ============================================================
\section{From Mobilization to Institutional Change}
%% ============================================================

The central question of this article is whether the creation of the
Department of Ucayali can be explained as gradual institutional change
or whether it represents a critical juncture produced by bottom-up
mobilization. To answer this question, the analysis evaluates the
sequence described in the previous sections in light of the theoretical
criteria discussed above. Following the literature on historical
institutionalism, the first step is to assess whether mechanisms of
incremental change are sufficient to account for the observed
transformation. Only if gradual explanations prove insufficient does
the hypothesis of a critical juncture become plausible
\autocite{streeck_introduction_2005,mahoney_explaining_2010}.

\subsection*{Ruling out gradual institutional change}

Gradualist accounts would interpret the creation of Ucayali as the
result of cumulative administrative adjustments within the existing
territorial framework. Mechanisms such as layering or conversion could
explain the addition of new administrative arrangements without altering
the basic structure of the state. However, the evidence suggests that
the transformation went beyond incremental adaptation.

First, the demand for a separate department did not emerge as a technical
administrative reform but as a political claim rooted in long-standing
conflicts over representation and resource distribution. For decades,
local actors had requested improvements within the existing departmental
structure, yet these requests produced only limited concessions. The
shift from demands for investment or administrative attention to the
demand for a new department represented a qualitative change in the
nature of the conflict.

Second, the institutional outcome involved the creation of a new
territorial unit with its own administrative authorities, budgetary
allocation, and political representation. Such a change cannot easily
be described as layering or conversion, since it altered the distribution
of authority between the center and the region rather than merely
modifying existing rules. In terms of the framework proposed by
\textcite{streeck_introduction_2005}, the transformation exceeded the scope
of gradual adjustment and instead redefined the institutional structure
itself.

Third, the timing of the reform suggests the importance of political
contingency. The decisive phase of the mobilization coincided with the
transition from military rule to civilian government, a period
characterized by uncertainty and reduced state capacity. Similar regional
demands had existed earlier, but they had not produced comparable results.
This pattern is consistent with the argument that institutional change
may depend not only on structural pressures but also on the presence of
a permissive political context \autocite{collier_shaping_1991}.

\subsection*{Antecedent tensions and permissive conditions}

The literature on critical junctures emphasizes the importance of
antecedent conditions that make multiple outcomes possible but do not
determine which one will occur \autocite{slater_informative_2010}. In the
present case, long-standing tensions between Pucallpa and the authorities
of Loreto created a situation in which demands for administrative
autonomy were both plausible and recurrent. Economic growth, demographic
expansion, and limited political representation generated pressures that
could be resolved through different institutional arrangements, including
decentralization, administrative reform, or the creation of a new
department.

These antecedent tensions alone, however, did not produce institutional
change. The decisive factor was the emergence of a permissive political
context during the late 1970s. The transition from military rule weakened
the capacity of the central government to impose decisions unilaterally
and increased the incentives to negotiate with regional actors. Under
these conditions, outcomes that had previously been unlikely became
politically feasible.

\subsection*{Productive mechanism: sustained mobilization}

A critical juncture requires not only permissive conditions but also a
productive mechanism linking contingency to institutional outcome
\autocite{capoccia_study_2007}. In the case of the Pucallpazo, this
mechanism was sustained collective mobilization at the subnational level.

The movement that emerged in Pucallpa was not a single protest but a
sequence of coordinated actions involving strikes, demonstrations,
negotiations, and public campaigns. The persistence of mobilization over
several years prevented the conflict from being resolved through
temporary concessions. Instead, it forced national authorities to
consider solutions that altered the territorial organization of the
state.

Importantly, the mobilization was decentralized and involved actors from
different sectors, including labor unions, business associations, civic
organizations, and local authorities. This broad coalition increased the
political cost of ignoring regional demands and made repression more
difficult during a period of national transition. The evidence suggests
that without this sustained pressure, the creation of a new department
would have remained unlikely.

\subsection*{Magnitude and durability of institutional change}

Another defining feature of critical junctures is the durability of
their consequences. Decisions taken during a brief period of contingency
must produce institutional effects that persist over time
\autocite{capoccia_study_2007}. The creation of the Department of
Ucayali meets this criterion.

The reform did not merely change policy; it altered the formal
territorial structure of the Peruvian state. The new department acquired
its own administrative authorities, budget, and representation within
national institutions. Once established, this arrangement proved durable
despite political changes at the national level. Subsequent governments
did not reverse the reform, and the new territorial structure became
part of the normal administrative framework of the country.

The scope of the transformation also distinguishes it from ordinary
policy concessions. Rather than addressing a specific grievance, the
reform redefined the relationship between the center and a peripheral
region. In this sense, the change was significant, swift relative to
previous decades of conflict, and encompassing in its institutional
effects.

\subsection*{Contingency and alternative outcomes}

Finally, the analysis must consider whether alternative outcomes were
realistically possible. A critical juncture implies that different
decisions could have led to different institutional paths
\autocite{capoccia_study_2007}. The historical record indicates that
several alternatives existed during the mobilization, including partial
concessions, administrative reforms without autonomy, or the postponement
of any decision.

Testimonies from participants and contemporaneous documents show that
the outcome remained uncertain even after the transition to civilian
rule. The central government faced competing pressures from different
regions and political actors, and the creation of a new department was
only one among several possible responses. The fact that this option was
ultimately adopted reflects the interaction between political uncertainty
and sustained regional mobilization rather than a predetermined
institutional trajectory.

\subsection*{A subnational critical juncture}

Taken together, the evidence supports the interpretation of the
Pucallpazo as a subnational critical juncture. Long-term center--periphery
tensions created the conditions for conflict, the transition to civilian
rule opened a window of opportunity, sustained mobilization acted as the
productive mechanism, and the resulting institutional change produced
durable effects on the territorial organization of the state.

This finding suggests that critical junctures need not originate at the
national level or be driven exclusively by elite negotiations. Under
certain conditions, collective action in peripheral regions can generate
institutional rupture and redefine authority relations within the state.

%% ============================================================
\section{Subnational Mobilization and Institutional Change\\
in Comparative Perspective}
%% ============================================================

The interpretation of the Pucallpazo as a subnational critical juncture
raises broader questions about the relationship between collective action,
territorial conflict, and institutional change in Latin America. Much of
the literature on democratization and institutional transformation in the
region has focused on national-level processes, particularly elite
negotiation, regime transitions, and constitutional reform
\autocite{huntington_tercera_2002,bunce_rethinking_2003}. While these approaches have
produced important insights, they often overlook the role of conflicts
originating in peripheral regions.

Research on contentious politics has shown that collective mobilization
can generate institutional effects even in highly centralized states,
particularly when protests are sustained and involve broad coalitions of
actors \autocite{tilly_social_2004,abers_inventing_2000}. However, most studies of social
movements emphasize policy change, participation, or the expansion of
rights, rather than transformations in the territorial organization of
the state.

Studies of subnational politics suggest that territorial dynamics often
shape institutional development in ways that cannot be understood solely
from the national level. \textcite{gibson_boundary_2005} argues that conflicts
over territorial authority are central to understanding the persistence
of unequal political arrangements within democratic regimes. Similarly,
\textcite{giraudy_politics_2010} shows that subnational political structures may
evolve independently from national political change, producing variation
in institutional outcomes across regions.

In Latin America, regional conflicts have frequently accompanied periods
of political transition or economic crisis, when central governments face
reduced capacity to control local actors. Comparative research on
institutional change suggests that these moments of uncertainty can create
opportunities for actors outside the national political center to
influence the direction of reform
\autocite{boone_political_2003,silva_challenging_2009}.

The Peruvian case illustrates this dynamic. During the late 1970s, the
transition from military rule to civilian government weakened the ability
of the central state to manage regional conflicts through routine
administrative mechanisms. At the same time, rapid demographic and
economic growth in the Amazon increased the salience of demands for
autonomy.

The Pucallpazo differs from more common trajectories in two respects.
First, the mobilization produced a coalition that included multiple social
sectors and maintained coordination over several years, increasing the
political cost of ignoring regional demands. Second, the events occurred
during a period of national uncertainty in which the central government
faced incentives to negotiate rather than repress. The combination of
these factors made possible an outcome that had previously appeared
unlikely: the creation of a new department and the redefinition of
territorial authority.

This case therefore suggests that institutional rupture can emerge from
the interaction between structural tensions, political contingency, and
sustained collective action at the subnational level. Recognizing this
possibility expands the scope of comparative research on democratization
and state formation in Latin America, highlighting the importance of
territorial conflicts as sources of institutional transformation.

%% ============================================================
\section{Conclusion}
%% ============================================================

This article has examined the Pucallpazo as an episode of institutional
change generated from the periphery of the Peruvian state. By
reconstructing the sequence that linked long-term center--periphery
tensions, the cycle of mobilization between 1978 and 1981, and the
creation of the Department of Ucayali, the analysis has argued that the
outcome cannot be adequately explained as gradual institutional
adaptation. Instead, the case is better understood as a subnational
critical juncture in which a relatively short period of political
uncertainty allowed collective actors to redefine the territorial
organization of the state.

The evidence shows that the reform did not result from the slow
accumulation of incremental adjustments, nor from a reform designed
entirely by national elites. For decades, regional demands had produced
only limited concessions without altering the basic structure of
authority. The creation of Ucayali occurred only when long-standing
structural tensions coincided with a permissive political context during
the transition from military to civilian rule, reducing the capacity of
the central government to manage conflict through routine mechanisms.
In this setting, decentralized mobilization transformed distributive
grievances into a demand for territorial autonomy and sustained pressure
long enough to make institutional reorganization a politically viable
solution.

The case contributes to the literature on institutional change by showing
that critical junctures can emerge outside the arenas most commonly
studied in historical institutionalism. Rather than resulting from
national crises or elite negotiation, the transformation analyzed here
originated in a peripheral region whose demographic growth and economic
importance were not matched by political representation. When this
imbalance became politically unsustainable, collective action did not
simply force policy concessions but produced an institutional rupture
that redefined authority relations between the center and the periphery.

The findings also speak to debates on democratization in Latin America.
The Pucallpazo illustrates a trajectory in which sustained mobilization
at the subnational level altered the territorial organization of the
state itself. This suggests that democratizing processes may unfold not
only through regime change but also through conflicts over territorial
authority, especially in highly centralized countries where regional
demands cannot easily be absorbed within existing institutional
arrangements.

More broadly, the case highlights the importance of peripheral regions
in the study of state formation and institutional change. Territorial
structures that appear stable may remain contingent when demographic
expansion, economic transformation, and limited political representation
generate pressures that existing institutions cannot accommodate. Under
such conditions, collective mobilization may produce outcomes that go
beyond reform and lead to durable reconfiguration of state authority.

Future research could examine whether similar subnational critical
junctures have occurred in other regions of Latin America, and under
what conditions regional mobilization succeeds in reshaping state
structures rather than merely influencing policy. By showing how
institutional rupture can originate from below, the Pucallpazo provides
a useful point of departure for analyzing the relationship between
territorial conflict, collective action, and democratic change in the
region.

%% ============================================================
%% References
%% ============================================================
\newpage
\printbibliography[title={References}]

%% ============================================================
%% Oral Sources (separate from bibliography per LARR style)
%% ============================================================
\subsection*{Oral Sources}

\begin{description}
  \item[Alfaro Venturo, Fernando.] 2022. Interview by [Author 3].
    Phone, April 2022.
  \item[Flores, Julio.] 2022. Interview by [Author 3].
    Phone, October 18, 2022.
  \item[Medina Ortiz, Eulogio.] 2022. Interview by [Author 3]. 2022.
  \item[Vásquez Valera, Manuel.] 2022. Interview by [Author 3].
    Phone, April 21, 2022.
  \item[Villa, Humberto.] 2022. Interview by [Author 3].
    Videoconference, October 15, 2022.
\end{description}

\end{document}
